% GENERATED FILE. Edit canonical lessons, then run npm run generate:books.
% canonical-source-hash: fnv1a64:a25a14c4a29e9c8c
% canonical-lessons: GU-C356-bandharan, GU-C356-svatantra, GU-C356-samaj, GU-C356-manya, GU-C356-jalavvu

\chapter{Rights and the Constitution}
\label{ch:gu-rights-and-the-constitution}
\hlchaptermodality{356}

\begin{chapteropening}
\textbf{By the end of this chapter:} \emph{I can talk about the rights a constitution protects.}

Complete the last lesson of chapter 356: I can talk about the rights a constitution protects.
\end{chapteropening}

\section[\texorpdfstring{\gu{બંધારણ} (bandhāraṇ)}{બંધારણ (bandhāraṇ)}]{\gu{બંધારણ} (bandhāraṇ) — a constitution}

\label{lesson:GU-C356-bandharan}



\begin{quote}
Before the new one: say the Gujarati for a grant, then the Gujarati for to collect money.
\end{quote}

\subsection*{You'll want to know: \gu{બંધારણ}}
\textbf{\gu{બંધારણ}} — \emph{bandhāraṇ} — \textquotedblleft{}a constitution\textquotedblright{}.

\begin{quote}
\gu{આપણું} \gu{બંધારણ} \gu{દરેક} \gu{નાગરિકને} \gu{સમાન} \gu{હક} \gu{આપે} \gu{છે}.
\end{quote}

\begin{quote}
\emph{āpṇũ bandhāraṇ darek nāgrikne samān hak āpe che.}
\end{quote}

\begin{quote}
Our constitution gives every citizen equal rights.
\end{quote}

A neuter noun.

\begin{cousinweb}
\textbf{\gu{બંધારણ}} (\emph{bandhāraṇ}) CDIAL lists Gujarati \emph{bandhāraṇ}, \textquotedblleft{}arrangement, construction\textquotedblright{}, under Sanskrit \emph{bandhakaraṇa}, \textquotedblleft{}binding\textquotedblright{} (CDIAL 9138); \textquotedblleft{}a constitution\textquotedblright{} is its modern sense, the frame a state is built on.
\end{cousinweb}

\begin{grammarlens}[title={What you've built}]
One more word for giving reasons and opinions.
\end{grammarlens}

\subsection*{Your turn}
\begin{itemize}
\raggedright
  \item \emph{Say it:} \emph{bandhāraṇ}
  \item \emph{Say it:} \emph{bandhāraṇ}, once more
  \item \emph{Say it:} \emph{āpṇũ bandhāraṇ darek nāgrikne samān hak āpe che.} — Our constitution gives every citizen equal rights.
  \item \emph{Recall:} say the Gujarati for a grant, then the Gujarati for to collect money, then say \emph{bandhāraṇ} again
\end{itemize}

\begin{tcolorbox}[breakable,colback=teal!4,colframe=teal!35!black,title={Before you move on}]
What does \gu{બંધારણ} mean? (\textquotedblleft{}A constitution\textquotedblright{}.) Say it once more.
\end{tcolorbox}
\section[\texorpdfstring{\gu{સ્વતંત્ર} (svatantra)}{સ્વતંત્ર (svatantra)}]{\gu{સ્વતંત્ર} (svatantra) — independent}

\label{lesson:GU-C356-svatantra}



\begin{quote}
Before the new one: say the Gujarati for to collect money, then the Gujarati for a constitution.
\end{quote}

\subsection*{You'll want to know: \gu{સ્વતંત્ર}}
\textbf{\gu{સ્વતંત્ર}} — \emph{svatantra} — \textquotedblleft{}independent\textquotedblright{}.

\begin{quote}
\gu{ન્યાયાધીશ} \gu{સ્વતંત્ર} \gu{હોવા} \gu{જોઈએ}.
\end{quote}

\begin{quote}
\emph{nyāyādhīś svatantra hovā joīe.}
\end{quote}

\begin{quote}
Judges should be independent.
\end{quote}

\begin{cousinweb}
\textbf{\gu{સ્વતંત્ર}} (\emph{svatantra}) is a tatsama, borrowed from Sanskrit \emph{svatantra}, \textquotedblleft{}self-dependence, freedom\textquotedblright{} (Monier-Williams): under one's own rule.
\end{cousinweb}

\begin{grammarlens}[title={What you've built}]
One more word for giving reasons and opinions.
\end{grammarlens}

\subsection*{Your turn}
\begin{itemize}
\raggedright
  \item \emph{Say it:} \emph{svatantra}
  \item \emph{Say it:} \emph{svatantra}, once more
  \item \emph{Say it:} \emph{nyāyādhīś svatantra hovā joīe.} — Judges should be independent.
  \item \emph{Recall:} say the Gujarati for to collect money, then the Gujarati for a constitution, then say \emph{svatantra} again
\end{itemize}

\begin{tcolorbox}[breakable,colback=teal!4,colframe=teal!35!black,title={Before you move on}]
What does \gu{સ્વતંત્ર} mean? (\textquotedblleft{}Independent\textquotedblright{}.) Say it once more.
\end{tcolorbox}
\section[\texorpdfstring{\gu{સમાજ} (samāj)}{સમાજ (samāj)}]{\gu{સમાજ} (samāj) — society}

\label{lesson:GU-C356-samaj}



\begin{quote}
Before the new one: say the Gujarati for a constitution, then the Gujarati for independent.
\end{quote}

\subsection*{You'll want to know: \gu{સમાજ}}
\textbf{\gu{સમાજ}} — \emph{samāj} — \textquotedblleft{}society\textquotedblright{}.

\begin{quote}
\gu{સારો} \gu{સમાજ} \gu{ગરીબ} \gu{માણસોની} \gu{કાળજી} \gu{લે} \gu{છે}.
\end{quote}

\begin{quote}
\emph{sāro samāj garīb māṇsonī kāḷjī le che.}
\end{quote}

\begin{quote}
A good society takes care of its poor.
\end{quote}

A masculine noun.

\begin{cousinweb}
\textbf{\gu{સમાજ}} (\emph{samāj}) is a tatsama, borrowed from Sanskrit \emph{samāja}, \textquotedblleft{}a meeting, an assembly\textquotedblright{} (Monier-Williams).
\end{cousinweb}

\begin{grammarlens}[title={What you've built}]
One more word for giving reasons and opinions.
\end{grammarlens}

\subsection*{Your turn}
\begin{itemize}
\raggedright
  \item \emph{Say it:} \emph{samāj}
  \item \emph{Say it:} \emph{samāj}, once more
  \item \emph{Say it:} \emph{sāro samāj garīb māṇsonī kāḷjī le che.} — A good society takes care of its poor.
  \item \emph{Recall:} say the Gujarati for a constitution, then the Gujarati for independent, then say \emph{samāj} again
\end{itemize}

\begin{tcolorbox}[breakable,colback=teal!4,colframe=teal!35!black,title={Before you move on}]
What does \gu{સમાજ} mean? (\textquotedblleft{}Society\textquotedblright{}.) Say it once more.
\end{tcolorbox}
\section[\texorpdfstring{\gu{માન્ય} (mānya)}{માન્ય (mānya)}]{\gu{માન્ય} (mānya) — valid; accepted}

\label{lesson:GU-C356-manya}



\begin{quote}
Before the new one: say the Gujarati for independent, then the Gujarati for society.
\end{quote}

\subsection*{You'll want to know: \gu{માન્ય}}
\textbf{\gu{માન્ય}} — \emph{mānya} — \textquotedblleft{}valid; accepted\textquotedblright{}.

\begin{quote}
\gu{આ} \gu{પરવાનો} \gu{ફક્ત} \gu{એક} \gu{વરસ} \gu{માટે} \gu{માન્ય} \gu{છે}.
\end{quote}

\begin{quote}
\emph{ā parvāno phakt ek varas māṭe mānya che.}
\end{quote}

\begin{quote}
This licence is valid for one year only.
\end{quote}

\begin{cousinweb}
\textbf{\gu{માન્ય}} (\emph{mānya}) is a tatsama, borrowed from Sanskrit \emph{mānya}, \textquotedblleft{}to be respected, to be honoured\textquotedblright{} (Monier-Williams); Gujarati uses it for what an office accepts as valid.
\end{cousinweb}

\begin{grammarlens}[title={What you've built}]
One more word for giving reasons and opinions.
\end{grammarlens}

\subsection*{Your turn}
\begin{itemize}
\raggedright
  \item \emph{Say it:} \emph{mānya}
  \item \emph{Say it:} \emph{mānya}, once more
  \item \emph{Say it:} \emph{ā parvāno phakt ek varas māṭe mānya che.} — This licence is valid for one year only.
  \item \emph{Recall:} say the Gujarati for independent, then the Gujarati for society, then say \emph{mānya} again
\end{itemize}

\begin{tcolorbox}[breakable,colback=teal!4,colframe=teal!35!black,title={Before you move on}]
What does \gu{માન્ય} mean? (\textquotedblleft{}Valid; accepted\textquotedblright{}.) Say it once more.
\end{tcolorbox}
\section[\texorpdfstring{\gu{જાળવવું} (jāḷavvũ)}{જાળવવું (jāḷavvũ)}]{\gu{જાળવવું} (jāḷavvũ) — to preserve}

\label{lesson:GU-C356-jalavvu}



\begin{quote}
Before the new one: say the Gujarati for society, then the Gujarati for valid.
\end{quote}

\subsection*{You'll want to know: \gu{જાળવવું}}
\textbf{\gu{જાળવવું}} — \emph{jāḷavvũ} — \textquotedblleft{}to preserve\textquotedblright{}.

\begin{quote}
\gu{આપણે} \gu{આપણા} \gu{હક} \gu{જાળવવા} \gu{જોઈએ}.
\end{quote}

\begin{quote}
\emph{āpaṇe āpṇā hak jāḷavvā joīe.}
\end{quote}

\begin{quote}
We must protect our rights.
\end{quote}

Past, with \gu{તેણે} (he or she …): \gu{તેણે} \gu{જૂનું} \gu{ઘર} \gu{જાળવ્યું}, \emph{teṇe jūnũ ghar jāḷavyũ}.

\begin{cousinweb}
\textbf{\gu{જાળવવું}} (\emph{jāḷavvũ}) CDIAL lists the Gujarati verb, \textquotedblleft{}to keep unharmed\textquotedblright{}, under a reconstructed \emph{yal}, \textquotedblleft{}endure, stay\textquotedblright{} (CDIAL 10430).
\end{cousinweb}

\begin{grammarlens}[title={What you've built}]
One more word for giving reasons and opinions.
\end{grammarlens}

\subsection*{Your turn}
\begin{itemize}
\raggedright
  \item \emph{Say it:} \emph{jāḷavvũ}
  \item \emph{Say it:} \emph{jāḷavvũ}, once more
  \item \emph{Say it:} \emph{āpaṇe āpṇā hak jāḷavvā joīe.} — We must protect our rights.
  \item \emph{Recall:} say the Gujarati for society, then the Gujarati for valid, then say \emph{jāḷavvũ} again
\end{itemize}

\begin{tcolorbox}[breakable,colback=teal!4,colframe=teal!35!black,title={Before you move on}]
What does \gu{જાળવવું} mean? (\textquotedblleft{}To preserve\textquotedblright{}.) Say it once more.
\end{tcolorbox}
